\chapter{Declaración sobre el uso de herramientas de IA}

Durante el desarrollo de este proyecto se han utilizado herramientas de inteligencia artificial como apoyo 
al trabajo realizado por los autores. En concreto, se ha recurrido a ChatGPT, Gemini y GitHub Copilot en 
distintas fases del proyecto, siempre como soporte y nunca como sustituto del criterio técnico o del trabajo 
de desarrollo, redacción y maquetación.

ChatGPT y Gemini se han empleado principalmente para:
\begin{itemize}
    \item resolver dudas sobre el uso de tecnologías como MongoDB y C\#;
    \item apoyar la redacción de partes de la memoria técnica, mejorando su claridad, estructura y redacción;
    \item proponer ideas para mejorar la lógica interna de la aplicación;
    \item generar ejemplos de peticiones y respuestas para entender mejor el funcionamiento de la API.
\end{itemize}

GitHub Copilot se ha utilizado durante el desarrollo del código y también en tareas de documentación. Su uso ha sido especialmente útil para:
\begin{itemize}
    \item autocompletar funciones en C\# y TypeScript;
    \item sugerir implementaciones en la lógica del backend y del frontend;
    \item acelerar la implementación del frontend, que no era un requisito inicial estricto del proyecto y se desarrolló como ampliación funcional;
    \item asistir en la redacción y revisión de la memoria técnica;
    \item facilitar el paso de una versión inicial redactada en Google Docs a su posterior maquetación en LaTeX.
\end{itemize}

Ejemplos de prompts utilizados:
\begin{itemize}
    \item ``Quiero que revises si la memoria cumple ahora mismo todos los requisitos que se piden en las instrucciones de entrega y en las cuestiones sobre los trabajos. Este es el README que hace de manual de despliegue y manual de usuario: README.md.''
    \item ``Quiero utilizar componentes de react-icons para los botones de la app, por ejemplo el de cerrar sesión o el más y menos de las cartas.''
    \item ``Quiero hacer un proyecto CRUD conectando C\# (.NET) con MongoDB. Necesito ideas de proyecto que tengan sentido; el proyecto lo realizamos entre dos personas.''
\end{itemize}

Algunas conversaciones con ChatGPT:
\begin{itemize}
    \item \url{https://chatgpt.com/share/69ee4684-5020-83eb-9fa2-355e1040c50c}
    \item \url{https://chatgpt.com/share/69ecf754-aeb0-8397-b190-0e320f49cae1}
\end{itemize}

En todo caso, el uso de estas herramientas se ha limitado a tareas de apoyo, revisión y sugerencia. Las decisiones 
finales de diseño, implementación, redacción y maquetación han sido tomadas por los autores del proyecto.
